% Section 16.6, from the summary table of lectures/L16/appendix/a_theory.md and from
% lectures/L16/README.md: the objectives, the evaluation questions and the next lesson.

\section{Sammanfattning}
\label{c16:sec:sammanfattning}

\begin{booktable}{@{}p{12.5em}L@{}}
  \thead{Begrepp} & \thead{Formel} \\
  \midrule
  Eulers formel & $e^{j\delta} = \cos\delta + j\sin\delta$ \\
  Eulerform & $z = \abs{z}e^{j\delta}$ \\
  Multiplikation & $z_1z_2 = \abs{z_1}\abs{z_2}\,\angle\,(\delta_1 + \delta_2)$ \\
  Division & $z_1/z_2 = (\abs{z_1}/\abs{z_2})\,\angle\,(\delta_1 - \delta_2)$ \\
  Från fasor till tidsdomän & $u(t) = \abs{U}\sin(\omega t + \delta)$ \\
\end{booktable}

\needspace{6\baselineskip}
\subsection*{Det här ska du kunna}
Efter det här kapitlet ska du:
\begin{itemize}
  \item Förstå Eulers formel.
  \item Känna till skillnaden mellan en komplex storhet och en fasor.
  \item Förstå hur komplexa tal kan användas vid vektorberäkningar.
\end{itemize}

\testadigsjalv
\begin{enumerate}
  \item Skriv $Z = 10\,\angle\,\frac{\pi}{6}$ på rektangulär form.
  \item Vad är värdet av $e^{j\pi}$? Vad innebär detta geometriskt?
\end{enumerate}

\needspace{6\baselineskip}
\subsection*{Nästa kapitel}
\Kapref{17} handlar om hur komplexa tal används för beräkningar med sinusformade signaler.
